%
% File: chap02.tex
% Author: Tengxiang Li
%
\chapter{Methodology of the Work}
\label{chap:04}
\setlength{\parskip}{1em}
%
The research and development of TinyCRM followed a scientific approach to entrepreneurship, utilizing the Lean Startup methodology to minimize market risk and maximize validated learning \cite{ries2011lean}, \cite{felin2024scientific}.

\section{Hypothesis Formulation and Testing}
The project began by identifying two core hypotheses. The value hypothesis assumed that micro-businesses in Kazakhstan would derive significant utility from an AI-driven tool that reduces manual data entry. The growth hypothesis assumed that a lightweight, mobile-first interface would drive adoption among solo entrepreneurs who work on the go. These were tested through customer development interviews and a quantitative survey of 15 target users.

\section{Quantitative Data Analysis}
The development utilized the Build-Measure-Learn feedback loop to iterate on product features \cite{ries2011lean}. A structured survey of 15 respondents provided empirical evidence for a strategic pivot. The survey results indicated that 86.7 percent of participants viewed manual data entry as a major bottleneck (rating 4 or 5 on a 5-point scale). Furthermore, 53.3 percent of users identified unexplained AI decisions as a deal-breaker, while 40 percent confirmed that their organizations prohibit the use of global generative AI for client data. These findings forced a shift from a fully autonomous system to a suggest and approve model, where the system provides recommendations along with explanations for the user to confirm.
